\section{Discussion}

The picture that emerges is consistent across every part of the argument. Closure drives
the economics through the $1/(1-C)$ leverage of (\ref{eq:leverage}); closure stops short of
one at the electronics wall; and the wall stands on two independent foundations, a supply
chain too deep to reproduce locally \cite{csis-semiconductor-supply-chain} and an embodied
energy some three orders of magnitude above ordinary structure
\cite{williams-ayres-heller-2002,gutowski-2009}. Because both foundations are physical
rather than a matter of engineering maturity, imported electronics are best treated as a
permanent line item. Modern seed-factory designs already do so, placing achievable closure
between about 70\% and \mbox{90-96\%} and declining to chase the rest
\cite{nasa-cp-2255-1980,guided-self-replicating-factory-2021,freitas-merkle-2004}.

\subsection{Limitations}
The energy side of the wall carries a measurement-basis caveat rather than a large
uncertainty. The chip figure is a finished-packaged-part basis; on a bare-wafer basis it
falls toward 1{,}800~kWh/kg and on an active-die basis it rises well above
\cite{murphy-2003,boyd-2012}, so the three-orders-of-magnitude gap is robust to the basis
but the exact multiple is not. The structural-metal end likewise uses terrestrial smelting
energy as a stand-in; in-situ regolith processing gives comparable single-digit-to-tens
figures but with its own plant mass and losses \cite{guerrero-zabel-2023}. The power-crossover
result is specific to one 97\%-closed lunar scenario and one replication model; the
qualitative conditional (local chip-making helps only when both closure and power are high)
is general, but the 4~MW threshold and the 29-to-17-year figures are scenario-dependent, in
the neighbourhood of the roughly 1.7~MW nominal seed power of the canonical study
\cite{nasa-cp-2255-1980}. Finally, the analysis is a mass-and-energy accounting, not a
manufacturability model: it says a lone seed cannot afford chips, not that a much larger,
richly powered industrial base never could.

\subsection{The narrow escape}
The one escape, having the factory make its own chips, is real but narrow. It helps only
for a factory that is both highly closed and richly powered, and it can actively hurt an
otherwise capable seed that is power-starved \cite{nasa-cp-2255-1980}. The practical
conclusion is not that self-replication fails, but that its design target is a factory
closed on nearly everything except chips, fed by imported vitamins, and sized so that its
power budget never tempts it into making the one part it cannot afford to make.
